\begin{textsample}{POS Dim 1 – 2007 – Score 24.00 – 2007/t000223.txt}  \label{ex:f1_pos_022}
Good afternoon, good evening, whatever.
We can go, jambo, guten Abend, bonsoir, but we can also ooh, ooh, ooh, ooh, ooh, ooh, ooh, ooh, ooh.
That is \textbf{the} call that chimpanzees make before they go to sleep in \textbf{the} evening.
You hear it going from one side of \textbf{the} valley to \textbf{the} other, from one group of nests to \textbf{the} next.
And I want to pick up with my talk this evening from where Zeray left off yesterday.
He was talking about this amazing, three-year-old Australopithecine child, Selam.
And we ' ve also been hearing about \textbf{the} history, \textbf{the} family tree, of mankind through DNA genetic profiling.
And it was a paleontologist, \textbf{the} late Louis Leakey, who actually set me on \textbf{the} path for studying chimpanzees.
And it was pretty extraordinary, way back then.
It ' s kind of commonplace now, but his argument was — because he ' d been \textbf{searching} for \textbf{the} fossilized remains of early humans in Africa.
And you can tell an awful lot about what those beings looked like from \textbf{the} \textbf{fossils}, from \textbf{the} shape of \textbf{the} muscle attachments, something about \textbf{the} way they lived from \textbf{the} various artifacts found with them.
But what about how they behaved?
That ' s what he wanted to know.
And of course, behavior doesn ' t fossilize.
He argued — and it ' s now a fairly common theory — that if we found behavior patterns similar or \textbf{the} same in our closest living \textbf{relatives}, \textbf{the} great apes, and humans today, then maybe those behaviors were present in \textbf{the} ape-like, human-like ancestor some seven million years ago.
And therefore, perhaps we had brought those characteristics with us from that ancient, ancient past.
Well, if you look in \textbf{textbooks} today that deal with human \textbf{evolution}, you very often find people \textbf{speculating} about how early humans may have behaved, based on \textbf{the} behavior of chimpanzees.
They are more like us than any other living creature, and we ' ve heard about that during this TED Conference.
So it remains for me to comment on \textbf{the} ways in which chimpanzees are so like us, in certain aspects of their behavior.
Every chimpanzee has his or her own personality.
Of course, I gave them names.
They can live to be 60 years or more, although we think most of them probably don ' t make it to 60 in \textbf{the} wild.
Mr.
Wurzel. \textbf{The} female has her first baby when she ' s 11 or 12.
Thereafter, she has one baby only every five or six years, a long period of childhood dependency when \textbf{the} child is nursing, sleeping with \textbf{the} mother at night, and riding on her back.
And we believe that this long period of childhood is important for chimpanzees, just as it is for us, in relation to learning.
As \textbf{the} brain becomes ever more complex during \textbf{evolution} in different forms of animals, so we find that \textbf{learning} plays an ever more important role in an individual ' s life history.
And young chimpanzees spend a lot of time watching what their elders do.
We know now that they ' re capable of imitating behaviors that they see.
And we believe that it ' s in this way that \textbf{the} different tool-using behaviors — that have now been seen in all \textbf{the} different chimpanzee populations studied in Africa — how these are passed from one generation to \textbf{the} next, through observation, imitation and practice, so that we can describe these tool-using behaviors as primitive culture.
Chimpanzees don ' t have a spoken language.
We ' ve talked about that.
They do have a very rich repertoire of postures and gestures, many of which are similar, or even \textbf{identical}, to ours and formed in \textbf{the} same context.
Greeting chimpanzees embracing.
They also kiss, hold hands, pat one another on \textbf{the} back.
And they swagger and they throw \textbf{rocks}.
In chimpanzee society, we find many, many examples of compassion, precursors to love and true altruism.
Unfortunately, they, like us, have a dark side to their nature.
They ' re capable of extreme brutality, even a kind of primitive war.
And these really aggressive behaviors, for \textbf{the} most part, are directed against individuals of \textbf{the} neighboring social group.
They are very territorially aggressive.
Chimpanzees, I believe, more than any other living creature, have helped us to understand that, after all, there is no sharp line between humans and \textbf{the} rest of \textbf{the} animal kingdom.
It ' s a very blurry line, and it ' s getting more blurry all \textbf{the} time as we make even more observations. \textbf{The} study that I began in 1960 is still continuing to this day.
And these chimpanzees, living their complex social lives in \textbf{the} wild, have helped — more than anything else — to make us realize we are part of, and not separated from, \textbf{the} amazing animals with whom we share \textbf{the} \textbf{planet}.
So it ' s pretty sad to find that chimpanzees, like so many other creatures around \textbf{the} world, are losing their \textbf{habitats}.
This is just one photograph from \textbf{the} air, and it shows you \textbf{the} forested highlands of Gombe.
And it was when I flew over \textbf{the} whole area, about 16 years ago, and realized that outside \textbf{the} park, this forest, which in 1960 had stretched almost unbroken along \textbf{the} \textbf{eastern} shore of Lake Tanganyika, which is where \textbf{the} tiny, 30-square-mile Gombe National Park lies, that a question came to my mind."How can we even try to save these \textbf{famous} chimpanzees, when \textbf{the} people living around \textbf{the} National Park are struggling to survive?"More people are living there than \textbf{the} land could possibly support. \textbf{The} numbers increased by refugees pouring in from Burundi and over \textbf{the} lake from Congo.
And very poor people — they couldn ' t afford to buy food from elsewhere.
This led to a program, which we call TACARE.
It ' s a very holistic way of improving \textbf{the} lives of \textbf{the} people living in \textbf{the} villages around \textbf{the} park.
It started small with 12 villages.
It ' s now in 24.
There isn ' t time to go into it, but it ' s including things like tree nurseries, methods of farming most suitable to this now very degraded, almost desert-like land up in these mountains.
Ways of controlling, preventing soil erosion.
Ways of reclaiming overused farmland, so that within two years they can again be productive.
Working to help \textbf{the} villagers obtain fresh water from wells.
Perhaps build some schoolrooms.
Most important of all, I believe, is working with small groups of women, providing them with opportunities for micro-credit loans.
And we ' ve got, as is \textbf{the} case around \textbf{the} world, about 95 percent of all loans returned.
Empowering women, working with education, providing scholarships for girls so they can finish secondary school, in \textbf{the} clear understanding that, all around \textbf{the} world, as women ' s education improves, family size drops.
We provide information about family planning and about HIV / AIDS.
And as a result of this program, something ' s happening for conservation.
What ' s happening for conservation is that \textbf{the} farmers living in these 24 villages, instead of looking on us as a bunch of white people coming to study a whole bunch of monkeys — and by \textbf{the} way, many of \textbf{the} staff are now Tanzanian — but when we began \textbf{the} TACARE program, it was a Tanzanian team going into \textbf{the} villages.
It was a Tanzanian team talking to \textbf{the} villagers, asking what they were interested in.
Were they interested in conservation?
Absolutely not.
They were interested in health ; they were interested in education.
And as time went on, and as their situation began to improve, they began to understand ever more about \textbf{the} need for conservation.
They began to understand that as \textbf{the} upper levels of \textbf{the} hills were denuded of trees, so you ' ve got this terrible soil erosion and mudslides.
Today, we are developing what we call \textbf{the} Greater Gombe Ecosystem.
This is an area way outside \textbf{the} National Park, stretching out into all these very degraded lands.
And as these villages have a better standard of life, they are actually agreeing to put between 10 percent and 20 percent of their land in \textbf{the} highlands aside, so that once again, as \textbf{the} trees grow back, \textbf{the} chimpanzees will have leafy corridors through which they can travel to interact — as they must for genetic viability — with other remnant groups outside \textbf{the} National Park.
So TACARE is a success.
We ' re replicating it in other parts of Africa, around other wilderness areas which are faced with extreme population pressure. \textbf{The} problems in Africa, however, as we ' ve been discussing for \textbf{the} whole of these first couple of days of TED, are major problems.
There is a great deal of poverty.
And when you get large numbers of people living in land that is not that fertile, particularly when you cut down trees, and you leave \textbf{the} soil open to \textbf{the} wind for erosion, as desperate populations cut down more and more trees, so that they can try and grow food for themselves and their families, what ' s going to happen?
Something ' s got to give.
And \textbf{the} other problems — in not only Africa, but \textbf{the} rest of \textbf{the} developing world and, indeed, everywhere — what are we doing to our \textbf{planet}?
You know, \textbf{the} \textbf{famous} scientist, E.
O.
Wilson said that if every person on this \textbf{planet} attains \textbf{the} standard of living of \textbf{the} average European or American, we need three new \textbf{planets}.
Today, they are saying four.
But we don ' t have them.
We ' ve got one.
And what ' s happened?
I mean, \textbf{the} question here is, here we are, arguably \textbf{the} most intelligent being that ' s ever walked \textbf{planet} Earth, with this extraordinary brain, capable of \textbf{the} kind of technology that is so well illustrated by these TED Conferences, and yet we ' re destroying \textbf{the} only home we have. \textbf{The} indigenous people around \textbf{the} world, before they made a major decision, used to sit around and ask themselves,"How does this decision affect our people seven generations ahead?"Today, major decisions — and I ' m not particularly talking about Africa here, but \textbf{the} developed world — major decisions involving millions of dollars, and millions of people, are often based on,"How will this affect \textbf{the} next shareholders ' meeting?"And these decisions affect Africa.
As I began traveling around Africa talking about \textbf{the} problems faced by chimpanzees and their vanishing forests, I realized more and more how so many of Africa ' s problems could be laid at \textbf{the} door of previous colonial exploitation.
So I began traveling outside Africa, talking in Europe, talking in \textbf{the} United States, going to Asia.
And everywhere there were these terrible problems.
And you know \textbf{the} kind I ' m talking about.
I ' m talking about pollution. \textbf{The} air that we breathe that often poisons us. \textbf{The} earth is poisoning our foods. \textbf{The} water — water is perhaps one of \textbf{the} most crucial issues that we ' re going to face in this century — and everywhere water is being \textbf{polluted} by agricultural, industrial and household chemicals that still are being sprayed around \textbf{the} world, seemingly with \textbf{the} inability to profit from past experience. \textbf{The} mangroves are being cut down ; \textbf{the} effects of things like \textbf{the} \textbf{tsunami} get worse.
We ' ve talked about \textbf{the} soil erosion.
We have \textbf{the} reckless burning of \textbf{fossil} fuels along with other \textbf{greenhouse} gasses, so called, leading to climate change.
Finally, all around \textbf{the} world, people have begun to believe that there is something going on very wrong with our climate.
All around \textbf{the} world climates are mixed up.
And it ' s \textbf{the} poor people who are affected worse.
It ' s Africa that already is affected.
In many parts of sub- Saharan Africa, \textbf{the} \textbf{droughts} are so much worse.
And when \textbf{the} rain does come, it so often leads to flooding and added distress, and \textbf{the} cycle of poverty and hunger and disease.
And \textbf{the} numbers of people living in an area that \textbf{the} land cannot support, who are too poor to buy food, who can ' t move away because \textbf{the} whole land is degraded.
And so you get desertification — creeping, creeping, creeping — as \textbf{the} last of \textbf{the} trees are cut down.
And this kind of thing is not just in Africa.
It ' s all over \textbf{the} world.
So it wasn ' t surprising to me that as I was traveling around \textbf{the} world I met so many young people who seemed to have lost hope.
We seem to have lost wisdom, \textbf{the} wisdom of \textbf{the} indigenous people.
I asked a question."Why?"Well, do you think there could be some kind of disconnect between this extraordinarily clever brain, \textbf{the} kind of brain that \textbf{the} TED technologies exemplify, and \textbf{the} human heart?
Talking about it in \textbf{the} non-scientific term, in terms of love and compassion.
Is there some disconnect?
And these young people, when I talk to them, basically they were either depressed or apathetic, or bitter and angry.
And they said more or less \textbf{the} same thing,"We feel this way because we feel you ' ve compromised our future and there ' s nothing we can do about it."We have compromised their future.
I ' ve got three little grandchildren, and every time I look at them and I think how we ' ve harmed this beautiful \textbf{planet} since I was their age, I feel this desperation.
And that led to this program we call Roots and Shoots, which began right here in Tanzania and has now spread to 97 countries around \textbf{the} world.
It ' s symbolic.
Roots make a firm foundation.
Shoots seem tiny ; to reach \textbf{the} sun they can break through a brick wall.
See \textbf{the} brick wall as all these problems we ' ve inflicted on \textbf{the} \textbf{planet}, environmental and social.
It ' s a message of hope.
Hundreds and thousands of young people around \textbf{the} world can break through and can make this a better world for all living things. \textbf{The} most important message of Roots and Shoots : every single one of us makes a difference, every single day.
We have a choice.
Every one of us in this room, we have a choice as to what kind of difference we want to make. \textbf{The} very poor have no choice.
It ' s up to us to change things so that \textbf{the} poor have choice as well. \textbf{The} Roots and Shoots groups all choose three projects.
It depends on how old they are, and which country, whether they ' re in a city or rural, as to what kinds of projects.
But basically, we have programs now from preschool right through university, with more and more adults starting their own Roots and Shoots groups.
And every group chooses, between them, three different kinds of project to make this a better world, recognizing that all these different problems are interconnected and impinge on each other.
So one of their projects will be to help their own human community.
And then, if they ' re able, they may raise money to help communities in other parts of \textbf{the} world.
One of their projects will be to help animals — not just wildlife, domestic animals as well.
And one of their projects will be to help \textbf{the} environment that we all share.
And woven throughout all of this is a message of learning to live in peace and harmony within ourselves, in our families, in our communities, between nations, between cultures, between religions and between us and \textbf{the} natural world.
We need \textbf{the} natural world.
We cannot go on destroying it at \textbf{the} rate we are.
We not do have more than this one \textbf{planet}.
Just picking one or two of \textbf{the} projects right here in Africa that \textbf{the} Roots and Shoots groups are doing, one or two projects only — in Tanzania, in Uganda, Kenya, South Africa, Congo-Brazzaville, Sierra Leone, Cameroon and other groups.
And as I say, it ' s in 97 countries around \textbf{the} world.
Of course, they ' re planting trees.
They ' re growing \textbf{organic} vegetables.
They ' re working in \textbf{the} refugee camps, with chickens and selling \textbf{the} eggs for a little amount of money, or just using them to feed their families, and feeling a sense of pride and empowerment, because they ' re no longer helpless and depending on others with their vegetables and their chickens.
It ' s being used in Uganda to give some psychological help to ex-child soldiers.
Doing projects like this is bringing them out of themselves.
Once again, they ' re useful members of society.
We have this program in prisons as well.
So, there ' s no time for more Roots and Shoots now.
But — oh, they ' re also working on HIV / AIDS.
That ' s a very important component of Roots and Shoots, with older kids talking to younger ones.
And unwanted pregnancies and things like that, which young people listen to better from other youth, rather than adults.
Hope.
That ' s \textbf{the} question I get asked as I ' m going around \textbf{the} world :"Jane, you ' ve seen so many terrible things, you ' ve seen your chimpanzees decrease in number from about one million, at \textbf{the} turn of \textbf{the} century, to no more than 150, 000 now, and \textbf{the} same with so many other animals.
Forests \textbf{disappearing}, deserts where once there was forest.
Do you really have hope?"Well, yes.
You can ' t come to a conference like TED and not have hope, can you?
And of course, there ' s hope.
One is this amazing human brain.
And I mean, think of \textbf{the} technologies.
And I ' ve just been so thrilled, finally, to come to people talking about compost latrines.
It ' s one of my hobbyhorses.
We just flush all this water down \textbf{the} lavatory, it ' s terrible.
And then talking about \textbf{renewable} energy — desperately important.
Do we care about \textbf{the} \textbf{planet} for our children?
How many of us have children or grandchildren, nieces, nephews?
Do we care about their future?
And if we care about their future, we, as \textbf{the} elite around \textbf{the} world, we can do something about it.
We can make choices as to how we live each day.
What we buy.
What we wear.
And choose to make these choices with \textbf{the} question, how will this affect \textbf{the} environment around me?
How will it affect \textbf{the} life of my child when he or she grows up?
Or my grandchild, or whatever it is.
So \textbf{the} human brain, coupled with \textbf{the} human heart, and we join hands around \textbf{the} world.
And that ' s what TED is helping so well with, and Google who help us, and Esri are helping us with mapping in Gombe National Park.
All of these technologies we can use.
Now let ' s link them, and it ' s beginning to happen, isn ' t it?
You ' ve heard about it this afternoon.
It ' s beginning to happen.
This change, this change.
To see change that we must have if we care about \textbf{the} future.
And \textbf{the} next reason for hope — nature is amazingly resilient.
You can take an area that ' s absolutely destroyed, with time and perhaps some help it can regenerate.
And an example is \textbf{the} TACARE program.
I told you, where a seemingly dead tree stump — if you stop hacking them for firewood, which you don ' t need to because you have wood lots, then in five years you can have a 30-foot tree.
And animals, almost on \textbf{the} brink of extinction, can be given a second chance.
That ' s my next book.
It ' s inspiring.
And it brings me to my last category of hope, and we ' ve heard about this so much in \textbf{the} last two days : this indomitable human spirit.
This determination of people, \textbf{the} resilience of \textbf{the} human spirit, So that people who you would think would be battered by poverty, or disease, or whatever, can pull themselves up out of it, sometimes with a helping hand, and take their part in society, and take their part in changing \textbf{the} world.
And just to think of one or two people out of Africa who are just really inspiring.
We could make a very long list, but obviously Nelson Mandela, emerging from 17 years of hard physical labor, 23 years of imprisonment, with this amazing ability to forgive, so that he could lead his nation out \textbf{the} evil regime of apartheid without a bloodbath.
Ken Saro-Wiwa, in Nigeria, who took on \textbf{the} giant \textbf{oil} companies, and although people around \textbf{the} world tried their best, was executed.
People like this are so inspirational.
People like this are \textbf{the} role models we need for young Africans.
And we need some environmental role models as well, and I ' ve been hearing some of them today.
So I ' m really grateful for this opportunity to share this message again, with everyone at TED.
And I hope that some of us can get together and talk about some of these things, especially \textbf{the} Roots and Shoots program.
And just a last word on that — \textbf{the} young woman who ' s running this entire conference center, I met her today.
She came up so excited, with her certificate.
She was [ in ] Roots and Shoots.
She was in \textbf{the} leadership in Dar es Salaam.
She said it ' s helped her to do what she ' s doing.
And it was very, very exciting for me to meet her and see just one example of how young people, when they are empowered, given \textbf{the} opportunity to take action, to make \textbf{the} world a better place, truly are our hope for tomorrow.
Thank you.

% matched lemmas: disappear, drought, eastern, evolution, famous, fossil, greenhouse, habitat, identical, learning, oil, organic, planet, pollute, relative, renewable, rock, search, speculate, textbook, the, tsunami
\end{textsample}
